\begin{table}[htbp]\centering\small
\caption{Clasificación de $V_n=V_0(1-hk)^n$ para $V_0,h,k>0$.}\label{tab:clasificacion}
\begin{tabular}{p{1.7cm}p{2.5cm}p{2.2cm}p{1.5cm}p{4.3cm}}
\toprule Régimen & Signo y monotonía & Magnitud & Límite & Estabilidad; sentido físico\\\midrule
$0<hk<1$ & Positivo; decrece & Disminuye & 0 & Asintótica; sí\\
$hk=1$ & Cero desde n=1 & Se anula & 0 & Asintótica; extinción artificial\\
$1<hk<2$ & Alternante & Disminuye & 0 & Asintótica; no\\
$hk=2$ & Alternante & Constante & No existe & Frontera; no\\
$hk>2$ & Alternante & Crece sin cota & No existe & Inestable; no\\
\bottomrule\end{tabular}
\par\smallskip\footnotesize Nota. En regímenes alternantes no hay monotonía de $V_n$. ``Frontera'' indica magnitud acotada sin amortiguación.
\end{table}